\subsection{Conditional Expressions and Predicates}
\label{Section 1.1.6}

The expressive power of the class of procedures that we can define at this
point is very limited, because we have no way to make tests and to perform
different operations depending on the result of a test. For instance, we
cannot define a procedure that computes the absolute value of a number by
testing whether the number is positive, negative, or zero and taking different
actions in the different cases according to the rule

$$
 |x| = \left\{ \begin{array}{r@{\quad \mathrm{if} \quad}l}
        x  &  x > 0, \\
	0  &  x = 0, \\
  \!\! -x  &  x < 0. \end{array} \right.
$$

This construct is called a \newterm{case analysis}.  Lisp notates it with a
special form called \code{cond}, which takes a list of cases and works down it
until one applies.  Python has no form of quite that shape, but it has a
\newterm{conditional expression}, which handles two cases, and a case analysis
of any length can be built by nesting one inside another:

\begin{codeBlock}
>>> def abs(x):
...     return x if x > 0 else (0 if x == 0 else -x)
\end{codeBlock}

\noindent
The general form of a conditional expression is

\begin{syntaxForm}
\meta{e}\textsubscript{1} if \meta{p}\textsubscript{1} else \meta{alternative}
\end{syntaxForm}

\noindent
and a case analysis of \meta{n} cases is written by putting another conditional
expression in the \meta{alternative} position, and so on down:

\begin{syntaxForm}
\meta{e}\textsubscript{1} if \meta{p}\textsubscript{1} else (
    \meta{e}\textsubscript{2} if \meta{p}\textsubscript{2} else (
        \dots
        \meta{e}\textsubscript{n} if \meta{p}\textsubscript{n} else \meta{e}\textsubscript{n+1}))
\end{syntaxForm}

\noindent
The pairs \meta{p} and \meta{e} are called \newterm{clauses}, and it must be
admitted that Python's way of writing them is the uglier of the two.  Lisp
lists the clauses one after another and they line up; Python nests them, and
the parentheses and indentation pile up on the right as the analysis grows.  We
will meet a more comfortable way to write this in a moment.

The first expression in each pair is a \newterm{predicate}---that is, an
expression whose value is interpreted as either true or false.\footnote{Both
languages accept any value where a predicate is expected, but they disagree
about which values count as false, and the disagreement is worth knowing early.
Scheme has two distinguished values, \code{\#t} and \code{\#f}, and treats
\code{\#f} as false and \textit{everything else}---including \code{0} and the
empty list---as true.  Python is more generous about what counts as false:
\code{False} and \code{None}, any number equal to zero, and any empty
collection (\code{''}, \code{[]}, \code{\{\}}, \code{()}) are all false, and
everything else is true.  The case to keep in mind is \code{0}, which is true
in Scheme and false in Python.  A type may decide the question for itself, and
we will see how in \link{Chapter 2}.}

Conditional expressions are evaluated as follows.  The predicate \( \langle{p_1}\rangle \) is
evaluated first.  If its value is false, then \( \langle{p_2}\rangle \) is evaluated.  If
\( \langle{p_2}\rangle \)'s value is also false, then \( \langle{p_3}\rangle \) is evaluated.  This process
continues until a predicate is found whose value is true, in which case the
interpreter returns the value of the corresponding \newterm{consequent
expression} \( \langle{e}\rangle \) of the clause as the value of the conditional expression.
If none of the \( \langle{p}\rangle \)'s is found to be true, the value of the \code{cond} is
undefined.

The word \newterm{predicate} is used for procedures that return true or false,
as well as for expressions that evaluate to true or false.  The absolute-value
procedure \code{abs} makes use of the primitive predicates \code{>}, \code{<},
and \code{==}.\footnote{\code{abs} also uses the ``minus'' operator \code{-},
which, when used with a single operand, as in \code{-x}, indicates
negation.  Note also that \code{abs} is built into Python already; we define it
here because the definition is the point, not the procedure.} These take two
numbers as arguments and test whether the first number is, respectively,
greater than, less than, or equal to the second number, returning true or false
accordingly.

The more comfortable way promised above is not an expression at all.  Python
has a case analysis in the form of a \newterm{statement}, written with
\code{if}, \code{elif} and \code{else}, and it lays its clauses out one after
another as Lisp does:

\begin{codeBlock}
>>> def abs(x):
...     if x > 0:
...         return x
...     elif x == 0:
...         return 0
...     else:
...         return -x
\end{codeBlock}

\noindent
This is how such a thing is ordinarily written in Python, and it is what we
will use from here on.  Note what has been given up to get it.  The conditional
expression \emph{has} a value, and can therefore appear anywhere a value can:
as an operand, as an argument, in the middle of a larger expression.  The
statement form has no value at all, and the branches must say \code{return}
explicitly to hand one back.  This is the same divide we met in
\link{Section 1.1.3}, and it will keep turning up: Python often gives us a
statement where Lisp gives us an expression, and the statement is more
comfortable to read and less free to compose.

\noindent
The absolute-value procedure could also be expressed in English as ``If
\( x \) is less than zero return \( -x; \) otherwise return \( x \),'' which in
Python is

\begin{codeBlock}
>>> def abs(x):
...     if x < 0:
...         return -x
...     else:
...         return x
\end{codeBlock}

\noindent
\code{else} names the clause taken whenever all the previous ones have been
bypassed.  Here there are precisely two cases, which is the shape the
conditional expression was made for; written that way it is

\begin{codeBlock}
>>> def abs(x):
...     return -x if x < 0 else x
\end{codeBlock}

\noindent
The general form of a conditional expression is

\begin{syntaxForm}
\meta{consequent} if \meta{predicate} else \meta{alternative}
\end{syntaxForm}

\noindent
Note that the predicate is written in the middle, between the two things it
chooses among, where Lisp writes it first.  To evaluate a conditional
expression, the interpreter starts by evaluating the \meta{predicate}.  If the
\meta{predicate} evaluates to a true value, the interpreter then evaluates the
\meta{consequent} and returns its value.  Otherwise it evaluates the
\meta{alternative} and returns its value.\footnote{A minor difference between
the two forms is that a branch of the statement form may be a sequence of
statements, evaluated in order.  In a conditional expression the
\meta{consequent} and \meta{alternative} must each be a single expression.}

In addition to primitive predicates such as \code{<}, \code{=}, and \code{>},
there are logical composition operations, which enable us to construct compound
predicates. The three most frequently used are these:

\begin{itemize}
\item
\meta{e}\textsubscript{1} \code{and} \dots\ \code{and} \meta{e}\textsubscript{n}

The interpreter evaluates the expressions \( \langle{e}\kern0.08em\rangle \) one at a time, in
left-to-right order.  If any \( \langle{e}\kern0.08em\rangle \) evaluates to false, the value of the
\code{and} expression is false, and the rest of the \( \langle{e}\kern0.08em\rangle \)'s are not
evaluated.  If all \( \langle{e}\kern0.08em\rangle \)'s evaluate to true values, the value of the
\code{and} expression is the value of the last one.

\item
\meta{e}\textsubscript{1} \code{or} \dots\ \code{or} \meta{e}\textsubscript{n}

The interpreter evaluates the expressions \( \langle{e}\kern0.08em\rangle \) one at a time, in
left-to-right order.  If any \( \langle{e}\kern0.08em\rangle \) evaluates to a true value, that value is
returned as the value of the \code{or} expression, and the rest of the
\( \langle{e}\kern0.08em\rangle \)'s are not evaluated.  If all \( \langle{e}\kern0.08em\rangle \)'s evaluate to false, the value
of the \code{or} expression is false.

\item
\code{not} \meta{e}

The value of a \code{not} expression is true when the expression \( \langle{e}\kern0.08em\rangle \)
evaluates to false, and false otherwise.
\end{itemize}

\noindent
Notice that \code{and} and \code{or} are built-in syntax, not procedures,
because the subexpressions are not necessarily all evaluated.  This is the
reason, promised in \link{Section 1.1.1}, that they are the two operators for
which no special method may be written: a procedure receives its arguments
already evaluated, and these two must be free not to evaluate one at all.
\code{not} has no such need, and is an ordinary operation.\footnote{One small
difference from Lisp: \code{and} and \code{or} return one of their operands, as
Lisp's do---\code{3 and 5} is \code{5}, and \code{0 or 'x'} is \code{'x'}---but
\code{not} always returns one of the two values \code{True} and \code{False},
never its operand.}

As an example of how these are used, the condition that a number \( x \) be in
the range \( 5 < x < 10 \) may be expressed as

\begin{codeBlock}
>>> x > 5 and x < 10
\end{codeBlock}

\noindent
though in this particular case Python offers something better.  Comparisons may
be \newterm{chained}, and written just as the mathematics is:

\begin{codeBlock}
>>> 5 < x < 10
\end{codeBlock}

\noindent
This is not merely an abbreviation.  \code{5 < x < 10} means \code{5 < x and x <
10} with the middle operand evaluated only once, which matters when it is
expensive to compute or has an effect, and it short-circuits in the same way:
if the first comparison is false the second is never made.  Chains may be of
any length, and may mix operators, as in \code{0 <= i < len(xs)}.

\noindent
The original goes on here to define a predicate for ``greater than or equal
to,'' which in Lisp must be built out of the two primitives it has.  Python
supplies it, and five others besides, so there is nothing to build.  It is
worth setting them out, since they are the whole stock of comparisons the
language has:

\begin{center}
\begin{tabular}{ll}
\code{x < y}   & \( x \) is less than \( y \) \\
\code{x <= y}  & \( x \) is less than or equal to \( y \) \\
\code{x > y}   & \( x \) is greater than \( y \) \\
\code{x >= y}  & \( x \) is greater than or equal to \( y \) \\
\code{x == y}  & \( x \) and \( y \) have the same value \\
\code{x != y}  & \( x \) and \( y \) do not have the same value \\
\end{tabular}
\end{center}

\noindent
Two of these deserve a remark.  The test for equality is \code{==}, two
characters, because the single \code{=} was spent on naming in
\link{Section 1.1.2}; Lisp, which names with \code{define}, had its \code{=}
free for the purpose.  And \code{!=} is the negation of \code{==}, so that
\code{x != y} and \code{not (x == y)} ask the same question.

None of these is confined to numbers.  \code{==} compares strings, lists and
anything else that knows how to compare itself, and the ordering comparisons
work on any type for which an order has been defined.  How a type comes to
define these for itself is the subject of \link{Chapter 2}; for the moment we
are dealing only with numbers, where they mean what they mean in arithmetic.

\begin{exerciseGroup}
\phantomsection\label{Exercise 1.1}
\exerciseHeading{Exercise 1.1:} Below is a sequence of expressions.
What is the result printed by the interpreter in response to each expression?
Assume that the sequence is to be evaluated in the order in which it is
presented.

\begin{codeBlock}
10
5 + 3 + 4
9 - 1
6 / 2
2 * 4 + (4 - 6)
a = 3
b = a + 1
a + b + a * b
a == b
b if b > a and b < a * b else a
\end{codeBlock}

\begin{codeBlock}
6 if a == 4 else (6 + 7 + a if b == 4 else 25)
\end{codeBlock}

\begin{codeBlock}
2 + (b if b > a else a)
\end{codeBlock}

\begin{codeBlock}
(a if a > b else (b if a < b else -1)) * (a + 1)
\end{codeBlock}

\phantomsection\label{Exercise 1.2}
\exerciseHeading{Exercise 1.2:} Write the following expression in
Python:

$$\frac{5 + 4 + (2 - (3 - (6 + \frac{4}{5})))}{3(6 - 2)(2 - 7)}.$$


\phantomsection\label{Exercise 1.3}
\exerciseHeading{Exercise 1.3:} Define a procedure that takes three
numbers as arguments and returns the sum of the squares of the two larger
numbers.

\phantomsection\label{Exercise 1.4}
\exerciseHeading{Exercise 1.4:} Observe that our model of
evaluation allows for compound expressions whose operators are themselves
computed.  The original writes the following procedure, in which the operator
position holds an \code{if} that chooses between \code{+} and \code{-}:

\begin{syntaxForm}
(define (a-plus-abs-b a b)
  ((if (> b 0) + -) a b))
\end{syntaxForm}

\noindent
Python cannot be written this way.  Its operators are syntax, as we saw in
\link{Section 1.1.1}, and \code{+} on its own is not an expression that can be
chosen between and then applied---it is a \code{SyntaxError}.  What Python
supplies instead is the \code{operator} module, which holds an ordinary
procedure for each operator: \code{operator.add} does what \code{+} does,
\code{operator.sub} what \code{-} does, \code{operator.mul}, \code{operator.lt}
and the rest likewise.  These are values, and so may be chosen between:

\begin{codeBlock}
>>> import operator
>>> def a_plus_abs_b(a, b):
...     return (operator.add if b > 0 else operator.sub)(a, b)
>>> a_plus_abs_b(3, -4)
~\textit{7}~
\end{codeBlock}

\noindent
Describe the behavior of this procedure.  Then consider what the detour through
\code{operator} has cost: the original needs no such module, because in Lisp
\code{+} was already an ordinary name for an ordinary procedure.  We have
recovered the ability, but not the uniformity that made it unremarkable.

\phantomsection\label{Exercise 1.5}
\exerciseHeading{Exercise 1.5:} Ben Bitdiddle has invented a test
to determine whether the interpreter he is faced with is using
applicative-order evaluation or normal-order evaluation.  He defines the
following two procedures:

\begin{codeBlock}
>>> def p():
...     return p()
>>> def test(x, y):
...     return 0 if x == 0 else y
\end{codeBlock}

Then he evaluates the expression

\begin{codeBlock}
>>> test(0, p())
\end{codeBlock}

What behavior will Ben observe with an interpreter that uses applicative-order
evaluation?  What behavior will he observe with an interpreter that uses
normal-order evaluation?  Explain your answer.  (Assume that the evaluation
rule for the conditional expression is the same whether the interpreter is
using normal or applicative order: The predicate expression is evaluated first,
and the result determines whether to evaluate the consequent or the alternative
expression.)
\end{exerciseGroup}

